\section{Discussion and Limitations}
\label{sec:discussion}

\subsection{Implications}
The proposed protocol encourages KT evaluation to move beyond aggregate AUC. Sparse-concept and calibration diagnostics can reveal whether a model is reliable where training evidence is limited. This is important for educational systems that use predicted probabilities for intervention or progression decisions.

\subsection{Scope Boundaries}
This paper is intentionally limited to evaluation diagnostics. It does not introduce a new KT model, a self-supervised learning objective, a graph neural network, a graph augmentation strategy, a learning path recommender, or a distillation method. Those directions require separate methodological papers and should be evaluated using the protocol introduced here.

\subsection{Limitations}
First, offline diagnostics do not prove improved learning outcomes in real classrooms. Second, limited cold-start definitions depend on dataset size and KC granularity. Third, some datasets may have incomplete Q-matrix provenance, especially when item--KC mappings are not strictly expert-tagged. Fourth, calibration findings may vary by dataset and baseline; therefore, ranking differences should be stated as observations under specific experimental conditions rather than universal claims.
